\newpage
\phantomsection
\label{prf:mcshane-equals-lebesgue}
\LRAProofFor{thm:mcshane-equals-lebesgue}
\begin{remark*}[Return]\hyperref[thm:mcshane-equals-lebesgue]{Return to Theorem}\end{remark*}
\begin{theorem*}[McShane Integrability Equals Lebesgue Integrability]
For real-valued functions on a compact interval,
\[
  \mathcal{McShane}[a,b]=\mathcal L[a,b],
\]
with agreement of integral values.
\end{theorem*}
\begin{proof}[Professional Standard Proof]
\LRAProofBodyStart
Show first that a Lebesgue-integrable function admits gauges for which every
McShane-fine sum is close to the Lebesgue integral. Approximate by simple
functions, control the exceptional level-set boundaries by small total
measure, and use free tags to force the same absolute-size bookkeeping as the
Lebesgue construction.

Conversely, show that McShane integrability implies Lebesgue integrability by
testing the free-tag sums against positive and negative parts. The freedom to
tag intervals from nearby points rules out the conditional cancellations that
HK permits, so the resulting indefinite integral is absolutely continuous and
has \(f\) as its derivative almost everywhere.
\end{proof}
\begin{proof}[Detailed Learning Proof]
\LRAProofBodyStart
The proof has two conceptual halves. Lebesgue integrability gives enough
measure control to make all free-tag gauge sums stable. McShane integrability
is strong enough to recover that same measure control, because the tags are
not tethered to the cells where cancellation would otherwise hide large
positive and negative oscillations.
\end{proof}
\begin{remark*}[Proof structure]
Reduce Lebesgue \(\Rightarrow\) McShane to simple-function approximation and
small exceptional sets. Reduce McShane \(\Rightarrow\) Lebesgue to absolute
continuity of the associated indefinite integral.
\end{remark*}
\begin{dependencies}
\begin{itemize}
  \item \hyperref[def:mcshane-fine-tagged-partition]{McShane-Fine Tagged Partition}.
  \item \hyperref[def:mcshane-integral]{McShane Integral}.
  \item \hyperref[thm:lebesgue-criterion-riemann-integrability]{Lebesgue Criterion for Riemann Integrability}.
\end{itemize}
\end{dependencies}
\clearpage
